% 各排版组合共用的虚构示例内容。
% 各主文件可在载入本文件前重定义下列组件，以形成差异化风格。
\providecommand{\ExampleCompanyOne}{%
  \resumebanner[ResumeSky][images/logo-stellar-tech.png]
    {星河智能科技有限公司}{AI 应用工程师}{2024.07 -- 至今}%
}
\providecommand{\ExampleCompanyTwo}{%
  \resumebanner[ResumePink][images/logo-cloud-lab.png]
    {云杉数据实验室}{算法工程师（实习）}{2023.07 -- 2024.03}%
}
\providecommand{\ExampleProjectOne}{%
  \resumeentry[ResumePrimary]{企业知识库问答平台}{RAG / 大语言模型应用}{2024.08 -- 2025.01}%
}
\providecommand{\ExampleProjectTwo}{%
  \resumeentry[ResumePrimary]{多智能体数据采集助手}{Agent / 浏览器自动化}{2024.03 -- 2024.07}%
}
\providecommand{\ExampleHighlight}[1]{\resumehighlight{#1}}

% 继承类默认前景叠层；裁掉占位照片文件自带的透明上边。
% photo-x-shift（正值向右）和 photo-y-shift（正值向上）可选，用于手动微调。
\resumeheader[images/placeholder-id-photo-transparent.png][photo-trim={0 0 0 168bp}]{示例姓名}{%
  \resumecontact{\faPhone}{联系电话：138-0000-0000}
  \resumecontact{\faEnvelope}{邮箱：\resumelink{mailto:hello@example.com}{hello@example.com}}\\
  \resumecontact{\faGlobe}{技术博客：\resumelink{https://example.com}{https://example.com} \resumehighlight{（10万+访问量）}}
  \\
  \resumecontact{\faGithub}{GitHub：\resumelink{https://github.com/your-name}{https://github.com/your-name}\hspace{0.35em}\resumehighlight[ResumeGold]{（\faStar\hspace{0.25em}100+ Star）}}
}

\resumesection{个人简介}
\begin{resumeitems}
  \item \textbf{个人概况：}关注大语言模型应用、数据分析与机器学习工程，具备从需求分析到服务部署的完整项目实践。
  \item \textbf{能力亮点：}熟悉 Python、PyTorch、RAG、Agent 与常用工程化工具。
\end{resumeitems}

\resumesection{工作经历}
\ExampleCompanyOne
\begin{resumeitems}
  \item 负责知识检索、智能问答与数据分析应用的算法设计、效果评估及服务化部署。
\end{resumeitems}

\ExampleCompanyTwo
\begin{resumeitems}
  \item 参与文本分类与信息抽取模型开发，完成数据清洗、误差分析和接口封装。
\end{resumeitems}

\resumesection{项目经历}
\ExampleProjectOne
\begin{resumeitems}
  \item \textbf{项目背景：}面向内部制度和产品手册构建可追溯的知识问答平台。
  \item \textbf{核心技术：}Python、FastAPI、混合检索、Rerank、Docker。
  \item \textbf{主要工作：}
    \begin{resumedetails}
      \item \textbf{数据管线：}\resumecircled{1}设计解析与索引流程；\resumecircled{2}统一多格式文档元数据。
      \item \textbf{检索优化：}\resumeparen{1}实现混合检索；\resumeparen{2}通过查询改写改善复杂问题质量。
      \item \textbf{效果评测：}\resumerightparen{1}建立评测集；\resumerightparen{2}准确率由 \ExampleHighlight{78\% 提升至 91\%}。
    \end{resumedetails}
\end{resumeitems}

\resumeprojectdivider
\ExampleProjectTwo
\begin{resumeitems}
  \item \textbf{项目背景：}将自然语言采集需求转换为可执行步骤，辅助整理公开网页数据。
  \item \textbf{主要工作：}设计规划、执行、页面分析与结果校验角色，使示例任务成功率达到 \ExampleHighlight{94\%}。
\end{resumeitems}

\resumesection{个人技能}
\begin{resumeitems}
  \item \textbf{编程与框架：}Python、SQL、PyTorch、FastAPI、Git、Docker。
  \item \textbf{大模型应用：}Prompt Engineering、RAG、Agent、Function Calling、模型评测。
\end{resumeitems}

\resumesection{教育背景}
\resumeentry[black]{海棠理工大学}{计算机科学与技术 \quad 本科}{2020.09 -- 2024.06}
\begin{resumeitems}
  \item \textbf{GPA：}3.7/4.0（专业前 10\%）。
  \item \textbf{荣誉奖项：}校级优秀学生奖学金、优秀毕业生。
  \item \textbf{主修课程：}数据结构、机器学习、数据库系统、自然语言处理。
\end{resumeitems}

\resumesection{学科竞赛}
\begin{resumeawards}
  \resumeaward{数据科学挑战赛全国一等奖}
  \resumeaward{高校算法应用示例赛二等奖}
  \resumeaward{人工智能创新竞赛银奖}
  \resumeaward{数学建模挑战赛一等奖}
\end{resumeawards}
